\chapter{刷题方法：从题目到可复用能力}

刷题和面试训练的核心不是“见过多少题”，而是能否把陌生题拆成熟悉结构。题目只是外壳，真正要识别的是数据怎么组织、状态怎么变化、答案依赖哪些历史信息、是否存在单调性、是否能淘汰候选、是否有重叠子问题。

面对一道题，第一步不是想模板，而是读约束。输入规模决定复杂度上限，值域决定能不能用计数数组或位图，有序性决定能不能双指针或二分，是否允许重复决定哈希表和去重逻辑，是否有负数决定滑动窗口和前缀和的选择。很多错误题解不是算法错，而是约束没看清。

第二步写暴力方法。暴力方法要明确枚举对象：枚举元素、下标对、子数组、子序列、路径、排列、状态，还是答案值。枚举对象一旦明确，复杂度就能估出来。比如两数之和枚举下标对是 \complexity{O(n^2)}；所有子数组是 \complexity{O(n^2)}；所有排列是 \complexity{O(n!)}；网格 BFS 扫所有格子是 \complexity{O(mn)}。

第三步找重复工作。重复工作有几种典型形态：反复查询某个历史值，适合哈希表；反复求区间和，适合前缀和；反复维护窗口内统计，适合滑动窗口；反复找当前最大或最小，适合堆或单调队列；反复搜索同一状态，适合记忆化或动态规划。

\begin{principlebox}{面试表达原则}
不要直接说“这题用哈希表”。更好的表达是：暴力会枚举所有下标对，时间是 \complexity{O(n^2)}；慢点在于每个数都要反复查找另一个数是否存在；我从左到右扫描，用哈希表保存已经看到的值到下标的映射；处理当前值时查 \code{target - nums[i]} 是否出现过；这样每个元素只插入和查询一次，平均时间 \complexity{O(n)}，空间 \complexity{O(n)}。
\end{principlebox}

第四步写不变量。不变量是优化算法正确性的核心。滑动窗口的不变量可能是“窗口内没有重复字符”；单调栈的不变量可能是“栈内下标对应的值单调递减”；BFS 的不变量是“当前层所有节点距离源点相同”；二分的不变量是“答案一定在当前左右边界之间”。没有不变量，代码只是碰巧通过样例。

第五步落到 \cpp。\cpp 代码要讲容器选择。\code{vector} 适合连续存储和下标访问；\code{unordered_map} 适合平均常数查询历史信息；\code{deque} 适合两端弹入弹出；\code{priority_queue} 适合动态极值但不能删除任意元素；\code{set} 适合有序前驱后继；图的邻接表通常用 \code{vector<vector<int>>}，稀疏图比邻接矩阵更节省空间。

仍然以两数之和为例。题目问的是两个不同下标的数能否凑成目标值，最直接的想法不是哈希表，而是枚举下标对。先固定左端 \code{left}，再枚举它右边的 \code{right}，检查二者之和是否等于 \code{target}。这段暴力代码覆盖了全部候选，因为任意合法答案都是某一对 \code{left < right} 的下标；它也天然避免同一个下标被用两次。

\begin{principlebox}{图示：两数之和从暴力到哈希}
\begin{tabular}{@{}p{0.19\linewidth}p{0.73\linewidth}@{}}
暴力枚举 & 固定 \code{left}，扫描 \code{right}，逐对检查 \code{nums[left] + nums[right]}。\\
重复工作 & 每个当前位置都在历史元素里重新寻找补数，历史查询没有被保存。\\
哈希状态 & \code{value_to_index} 保存已经见过的“值 -> 下标”，只表示当前下标之前的历史。\\
循环顺序 & 先查 \code{target - nums[index]}，再插入 \code{nums[index]}，保证不会使用同一个下标两次。\\
\end{tabular}
\end{principlebox}

\begin{lstlisting}[language=C++,caption={LeetCode 1 两数之和：先写完整暴力基线}]
std::vector<int> two_sum_bruteforce(const std::vector<int>& nums, int target)
{
    for (int left = 0; left < static_cast<int>(nums.size()); ++left) {
        for (int right = left + 1; right < static_cast<int>(nums.size()); ++right) {
            const long long sum =
                static_cast<long long>(nums[left]) + nums[right];
            if (sum == target) {
                return {left, right};
            }
        }
    }
    return {};
}
\end{lstlisting}

暴力版本慢在“查找补数”这一步。固定当前数 \code{nums[right]} 时，我们真正需要知道的是此前是否出现过 \code{target - nums[right]}，不需要把所有历史下标重新扫一遍。于是优化不是凭空说“用哈希表”，而是把历史值组织成值到下标的映射：扫描到当前下标时，先查询补数是否已经在历史表中；如果在，答案就是历史下标和当前下标；如果不在，再把当前值放进历史表，供后面的元素查询。

\begin{lstlisting}[language=C++,caption={LeetCode 1 两数之和：哈希表把补数查找变成历史查询}]
std::vector<int> two_sum_hash_table(const std::vector<int>& nums, int target)
{
    std::unordered_map<int, int> value_to_index;
    value_to_index.reserve(nums.size());

    for (int index = 0; index < static_cast<int>(nums.size()); ++index) {
        const int need = target - nums[index];
        const auto found = value_to_index.find(need);
        if (found != value_to_index.end()) {
            return {found->second, index};
        }
        value_to_index.emplace(nums[index], index);
    }
    return {};
}
\end{lstlisting}

这段优化代码的不变量是：处理下标 \code{index} 前，哈希表只保存 \code{[0, index)} 中已经出现过的值。先查再插入，正是为了保持“历史”这个语义，避免当前元素和自己配对。参数用 \code{const std::vector<int>&}，避免复制输入数组。\code{reserve} 不是必须，但能减少扩容和 rehash。返回空数组表示没有答案，真实力扣题如果保证有答案，也可以直接返回。这样讲完以后，读者能看到完整链条：枚举下标对可以保证正确，瓶颈是重复扫描历史，哈希表只负责把历史补数查询降到平均常数。

\section{训练路线：从基础题到综合能力}

学习路线不是倒计时。有人一个月集中准备，有人一边工作一边长期训练，也有人先补 \cpp 和数据结构基础，再进入高强度题目练习。因此本章只给出能力阶段和训练节奏，不把整本书绑定到固定月份。你可以把每个阶段压缩成一周，也可以展开成一个月；真正不能省略的是每个阶段要掌握的知识、题型、代码能力和复盘方法。

第一阶段建立基础语言和复杂度意识。目标是能读懂输入规模，估计可接受复杂度，写出暴力解，并说清楚暴力解慢在哪里。这个阶段学习数组、字符串、\code{vector}、\code{string}、排序、二分库函数、下标边界、引用传参和整数溢出。推荐题包括 LeetCode 1 两数之和、LeetCode 283 移动零、LeetCode 26 删除有序数组重复项、LeetCode 88 合并两个有序数组、LeetCode 53 最大子数组和、LeetCode 344 反转字符串和 LeetCode 125 有效回文。做题时不要急着背模板，先把枚举对象写清楚：枚举元素、下标对、区间，还是答案边界。

第二阶段学习线性结构和窗口类算法。双指针、滑动窗口、前缀和、单调栈和单调队列看起来像技巧，本质都是“线性序列上的状态维护”。窗口题要写出左右边界的含义，说明什么时候扩张、什么时候收缩；单调结构题要说明哪些候选被淘汰，为什么淘汰后不会丢答案。推荐题包括 LeetCode 11 盛最多水的容器、LeetCode 15 三数之和、LeetCode 209 长度最小的子数组、LeetCode 3 无重复字符的最长子串、LeetCode 438 找到字符串中所有字母异位词、LeetCode 76 最小覆盖子串、LeetCode 739 每日温度、LeetCode 84 柱状图最大矩形和 LeetCode 239 滑动窗口最大值。

第三阶段把哈希、堆、树和图串起来。哈希表负责历史查询，堆负责动态极值，树负责层次和有序关系，图负责任意状态转移。这个阶段要形成建模习惯：节点是什么，边是什么，状态保存在哪里，访问标记何时更新，答案在入队、出队、入栈、出栈还是回溯时产生。推荐题包括 LeetCode 49 字母异位词分组、LeetCode 128 最长连续序列、LeetCode 560 和为 K 的子数组、LeetCode 347 前 K 个高频元素、LeetCode 215 数组中的第 K 个最大元素、LeetCode 102 二叉树层序遍历、LeetCode 98 验证二叉搜索树、LeetCode 236 二叉树的最近公共祖先、LeetCode 200 岛屿数量、LeetCode 994 腐烂的橘子、LeetCode 207 课程表和 LeetCode 210 课程表 II。

第四阶段进入搜索、回溯和动态规划。回溯题要写清选择列表、路径状态、终止条件和撤销动作；动态规划题要写清状态定义、转移来源、初始化、遍历顺序和空间优化条件。不要先背方程，先问一个状态是否足够描述未来。推荐题包括 LeetCode 46 全排列、LeetCode 78 子集、LeetCode 39 组合总和、LeetCode 22 括号生成、LeetCode 79 单词搜索、LeetCode 70 爬楼梯、LeetCode 198 打家劫舍、LeetCode 322 零钱兑换、LeetCode 300 最长递增子序列、LeetCode 1143 最长公共子序列、LeetCode 72 编辑距离、LeetCode 416 分割等和子集和 LeetCode 494 目标和。

第五阶段处理贪心、并查集和综合题。贪心题的关键是证明局部选择为什么能推出全局最优；并查集题的关键是把动态连通关系建模成集合合并；综合题则常常把哈希、排序、堆、图和 DP 混在一起。推荐题包括 LeetCode 55 跳跃游戏、LeetCode 134 加油站、LeetCode 56 合并区间、LeetCode 684 冗余连接、LeetCode 721 账户合并、LeetCode 208 实现 Trie、LeetCode 1319 连通网络的操作次数、LeetCode 127 单词接龙和 LeetCode 295 数据流中位数。这个阶段不要只追求 AC，要训练面试表达：暴力解是什么，瓶颈在哪里，优化依赖哪条性质，代码如何保证边界正确。

如果需要集中训练，可以把这些阶段压成十二周：基础复杂度与数组字符串，双指针与窗口，哈希与前缀和，栈队列堆与单调结构，二分查找，链表与树，图和拓扑排序，回溯，一维动态规划，二维动态规划和背包，贪心与并查集，综合模拟与查漏补缺。这只是节奏表，不是本书的边界。书的目标是把知识完整讲清楚，训练周期由你的基础、时间和目标岗位决定。

\begin{principlebox}{阶段交付}
每个阶段至少留下四类材料：一组代表题的暴力解和优化解，一份数据结构选择说明，一份 \cpp 容器和边界错误清单，一份一周后能独立重写的回访题表。刷题报告不是形式，它是把题解转化成长期能力的工具。
\end{principlebox}

\section{每日训练闭环：选题、限时、复盘和回访}

真正稳定的刷题节奏不是每天随机点开几道题，而是把每道题放进同一个闭环。一次训练至少包含四步：先选题，限定当前要训练的能力；再限时推导，强迫自己写出暴力枚举对象和优化依据；随后提交代码并记录失败原因；最后安排回访，确认一周后还能独立重写。

每天不要同时训练太多新模型。更稳的组合是“一道新题、一题同模型变形、一题回访题”。新题用来扩展边界，变形题用来确认模型能迁移，回访题用来发现自己到底记住了推理还是只记住了答案。若今天学习前缀和，就不要同时混入全新的图论和动态规划；可以选择 LeetCode 303 区域和检索、LeetCode 560 和为 K 的子数组、LeetCode 525 连续数组这一组，因为它们都围绕“前缀状态保存什么”展开。

\begin{principlebox}{一题一卡}
每道题复盘只保留六个字段：暴力枚举对象、重复工作、优化依赖的性质、核心数据结构、不变量、失败样例。题解代码可以很长，复盘卡必须短；短到一周后重新打开时，能立刻知道自己应该重建哪条推理链。
\end{principlebox}

\begin{lstlisting}[numbers=none,caption={每日刷题记录的最小字段}]
date
problem
topic
bruteforce_object
bottleneck
invariant
container_choice
failed_case
revisit_date
\end{lstlisting}

以 \code{560. 和为 K 的子数组} 为例，暴力枚举对象是所有闭区间 \code{[left,right]}，瓶颈是每个左端都要重复累加后缀。优化并不是“看到子数组就滑动窗口”：数组允许负数，窗口和不具备随右端扩张单调增加、随左端收缩单调减少的性质。正确的优化条件是前缀和等式：若当前前缀为 \code{prefix}，此前出现过 \code{prefix-k}，那么这些出现位置到当前下标之间的子数组和就是 \code{k}。因此哈希表保存的不是元素值，而是“历史前缀和出现次数”。

这类题的回访标准也要具体。第一次通过后，不看题解重写一次暴力版和前缀哈希版；第二天改输入为包含负数和零的数组，解释为什么正数滑动窗口不适用；一周后再做 \code{525. 连续数组} 或 \code{974. 和可被 K 整除的子数组}，检查自己能否把“保存前缀状态”迁移到差值和余数。

\section{学习方法和题解复盘的深层要求}

刷题教材不能只给最终代码。读者真正需要的是从题目约束到算法选择的推理链：输入规模允许什么复杂度，数据是否有序，值域能否压缩，状态是否会重复，局部选择能否证明，容器操作成本是否符合题目热路径。每道题都要先写暴力解，因为暴力解暴露了最朴素的状态空间；然后再指出哪个查询、哪个枚举、哪个重复子问题最贵。优化不是把模板贴上去，而是把最贵的那一步替换为更合适的数据结构或更强的不变量。

高质量复盘至少包含三层。第一层是题面复述，用自己的话说明输入、输出和约束，尤其要把隐藏条件翻译成算法语言，例如正数数组带来窗口单调性，排序数组带来二分和双指针，树节点唯一性决定能否用值建索引。第二层是复杂度预算，把 n、m、值域、字符串长度、图边数分开，不要只写一个 n。第三层是正确性说明，用不变量、交换论证、归纳或最短路层次证明解释为什么不会漏答案。

题解里的代码要服务于解释，而不是炫技。变量名应该表达状态含义，例如 left 表示窗口左边界，farthest 表示当前可达最远位置，indegree 表示尚未满足的依赖数量。复杂代码要先拆出检查函数、松弛函数或合并函数。这样做不是为了形式，而是让读者能在面试中逐句讲清楚代码对应的数学状态。代码写完以后，至少用一个最小用例、一个边界用例、一个反例用例和一个规模用例手算。

三个月计划只能作为节奏，不应该成为内容边界。真正的学习路径是循环式的：第一轮建立题型地图，第二轮补齐证明和边界，第三轮把相似题横向比较，第四轮把代码改成工程上也能接受的版本。每一轮都要记录失败原因。失败原因不是笼统的不会，而是要拆成读错题、复杂度预算错误、状态定义不完整、边界初始化错误、容器 API 误用、证明缺失或代码实现失误。

面试表达要像讲工程方案。先说暴力方案，说明它为什么慢；再说关键观察，说明它让哪一步变便宜；然后给出数据结构职责，说明保存什么、查询什么、删除什么；最后给出不变量和复杂度。这个顺序能防止答案变成模板堆砌。面试官追问时，优先回到约束和不变量，而不是机械背另一段代码。

实验是把知识落地的关键。数组题打印指针区间，窗口题打印计数表和 formed，二分题打印 left、mid、right 和谓词结果，图题打印队列层次，DP 题打印小表格，堆题打印堆顶和过期状态。只要一次实验能解释一个常见错误，它就不是额外负担，而是把题解变成长期能力的工具。

\begin{principlebox}{专题复盘要求}
读完本节后，不要只记结论。请至少回到 LeetCode 53 最大子数组和、LeetCode 560 和为 K 的子数组、LeetCode 875 爱吃香蕉、LeetCode 312 戳气球这类代表题，把每个结论还原成一个可验证的问题：它解决了哪一步重复工作，依赖哪个题目条件，使用哪个容器或状态，边界条件如何破坏错误写法。
\end{principlebox}


\chapter{复杂度、暴力法和优化来源}

复杂度是算法设计的预算。输入规模如果是 \code{n <= 100}，\complexity{O(n^3)} 可能可以接受；如果是 \code{n <= 2 \times 10^5}，通常要接近 \complexity{O(n \log n)} 或 \complexity{O(n)}；如果网格规模是 \code{m * n <= 10^5}，BFS 或 DFS 扫一遍通常合理；如果状态空间是指数级，就要考虑剪枝、记忆化或动态规划。

暴力法的作用是把问题完整展开。两数之和暴力枚举所有下标对，最长无重复子串暴力枚举所有起点和终点，岛屿数量暴力从每个未访问陆地扩展连通块。暴力方法慢，但它通常是正确性的起点。

\begin{lstlisting}[language=C++,caption={两数之和暴力解}]
std::vector<int> two_sum_bruteforce(const std::vector<int>& nums, int target)
{
    for (int left = 0; left < static_cast<int>(nums.size()); ++left) {
        for (int right = left + 1; right < static_cast<int>(nums.size()); ++right) {
            if (nums[left] + nums[right] == target) {
                return {left, right};
            }
        }
    }
    return {};
}
\end{lstlisting}

这段代码枚举的是下标对，数量约为 \code{n * (n - 1) / 2}，时间复杂度是 \complexity{O(n^2)}。它慢在每个数都要重新扫描后面的数。优化思路不是凭空出现，而是把“扫描后面所有数”改成“快速查询某个补数是否出现过”。这就是哈希表。

常见优化来源有六类。第一类是缓存重复计算，例如前缀和、记忆化搜索和动态规划。第二类是单调性，例如二分、双指针和某些滑动窗口。第三类是局部状态可维护，例如字符计数、窗口和、不同元素个数。第四类是候选淘汰，例如单调栈和单调队列。第五类是数据结构查询能力，例如哈希表、堆、平衡树、并查集。第六类是改变建模方式，例如把网格问题建成图，把区间问题排序，把字符串问题转为频次。

写完暴力法后，应该固定问几个问题：我枚举的对象是什么？有没有重复计算同一段信息？是否只需要历史信息的一小部分？是否存在有序或单调？是否可以把区间查询变成前缀差？是否可以用哈希把线性查询变成平均常数？是否可以把指数搜索合并成状态 DP？

面试时要把这些问题讲出来。一个好答案通常不是从“我用某某模板”开始，而是从暴力法、瓶颈和优化依据开始。

